\documentclass[10pt,a4paper,twoside,openany]{book}

% ── Encoding and Language ──────────────────────
\usepackage[utf8]{inputenc}
\usepackage[T1]{fontenc}
\usepackage[english]{babel}

% ── Mathematics and Physics ────────────────────
\usepackage{amsmath,amssymb,amsthm}
\usepackage{mathtools}
\usepackage{physics}
\usepackage{siunitx}

% ── Font (Arial-like, Greek-capable) ───────────
\usepackage{fontspec}
\setmainfont{Noto Sans}
\setsansfont{Noto Sans}
\setmonofont{Noto Sans Mono}

% ── Sans-serif for text, serif for math ────────
\usepackage{unicode-math}
\setmathfont{Latin Modern Math}
\setmathrm{Latin Modern Math}

% ── Page Geometry ──────────────────────────────
\usepackage[margin=2.5cm]{geometry}

% ── Graphics ───────────────────────────────────
\usepackage{graphicx}
\usepackage{tikz}
\usepackage{pgfplots}
\pgfplotsset{compat=1.18}

% ── Theorem Environments (shared counter) ──────
\theoremstyle{plain}
\newtheorem{theorem}{Theorem}[section]
\newtheorem{lemma}[theorem]{Lemma}
\newtheorem{proposition}[theorem]{Proposition}
\newtheorem{corollary}[theorem]{Corollary}

\theoremstyle{definition}
\newtheorem{definition}[theorem]{Definition}
\newtheorem{example}[theorem]{Example}
\newtheorem{exercise}[theorem]{Exercise}

\theoremstyle{remark}
\newtheorem{remark}[theorem]{Remark}
\newtheorem{note}[theorem]{Note}

% ── Equation numbering within sections ─────────
\numberwithin{equation}{section}

% ── Hyperlinks (plain black) ───────────────────
\usepackage[colorlinks=true,
            linkcolor=black,
            citecolor=black,
            urlcolor=black]{hyperref}

% ── Chapter Formatting (no "Chapter X") ────────
\usepackage{titlesec}
\titleformat{\chapter}
  {\normalfont\large\bfseries\raggedright}
  {\thechapter}{1em}{}
\titlespacing*{\chapter}{0pt}{-20pt}{4pt}

\titleformat{\section}
  {\normalfont\normalsize\bfseries\raggedright}
  {\thesection}{1em}{}
\titlespacing*{\section}{0pt}{4pt}{2pt}

\titleformat{\subsection}
  {\normalfont\normalsize\itshape\raggedright}
  {\thesubsection}{1em}{}
\titlespacing*{\subsection}{0pt}{3pt}{1pt}

% ── Tight Spacing ──────────────────────────────
\setlength{\parskip}{0.0em}
\setlength{\parindent}{12pt}
\setlength{\abovedisplayskip}{4pt}
\setlength{\belowdisplayskip}{4pt}
\setlength{\abovedisplayshortskip}{2pt}
\setlength{\belowdisplayshortskip}{2pt}

% ── List Spacing ───────────────────────────────
\usepackage{enumitem}
\setlist{nosep,leftmargin=*,itemsep=0pt,topsep=0pt,parsep=0pt}

% ── Plain Page Style ───────────────────────────
\pagestyle{plain}

% ── Title ──────────────────────────────────────
\title{\normalsize\bfseries A Concise Introduction to\\[0.2em] Mathematical Physics}
\author{Your Name}
\date{\today}

% ════════════════════════════════════════════════
\begin{document}
% ════════════════════════════════════════════════

\frontmatter
\maketitle
\tableofcontents

\chapter*{Preface}
\addcontentsline{toc}{chapter}{Preface}
TODO

\mainmatter

\chapter{Απόκλιση}\label{chapter_divergence}

Ως \textit{απόκλιση} διανυσματικού πεδίου εννοείται το βαθμωτό μέγεθος
που εκφράζει τον τοπικό ρυθμό μεταβολής του όγκου 
κατά μήκος των δυναμικών γραμμών του πεδίου.
Ακριβέστερα, σε κάθε σημείο του χώρου,
η απόκλιση ισούται με τον σχετικό ρυθμό ογκομετρικής μεταβολής
κατά μήκος της ροής του πεδίου,
δηλαδή τον οριακό ρυθμό αυξομείωσης ενός ρέοντος χωρίου ανά μονάδα όγκου
καθώς αυτό συρρικνώνεται αυθαίρετα γύρω από το υπό μελέτη σημείο.

% \operatorname{div}\xi(p)
% =
\begin{equation}
\lim_{\substack{V \to \{p\}}}
\frac{1}{\operatorname{vol} (V)}
\left.\frac{d}{dt}\right|_{t=0}
\hspace{-10pt}
\operatorname{vol} \bigl(\phi_t(V)\bigr)
\end{equation}

\noindent
Η απόκλιση είναι θετική ή αρνητική
αναλόγως του εάν το απειροστό στοιχείο διαστέλλεται ή συστέλλεται
καθώς ρέει,
και μηδενική εάν ο όγκος του παραμένει αναλλοίωτος
(τουλάχιστον σε πρώτο βαθμό προσέγγισης).
Η απόλυτη τιμή της απόκλισης εκφράζει την τοπική ένταση
της μεταβολής του παρασυρόμενου όγκου.

Ο όγκος ενός ρέοντος χωρίου μπορεί να διαστέλλεται
επειδή οι δυναμικές γραμμές του πεδίου τείνουν να αποκλίνουν γεωμετρικά
ή η ταχύτητα της ροής αυξάνεται κατά μήκος τους
(ή και τα δύο)· παρομοίως, η συστολή των παρασυρόμενων χωρίων
υποδηλώνει ότι υπερισχύει συνολικά κάποια σύγκλιση ή επιβράδυνση της ροής
(ή κάποιος συνδυασμός και των δύο).
Κατά συνέπεια, υποθέτοντας προς στιγμή πεδίο σταθερών ταχυτήτων,
η απόκλιση ποσοτικοποιεί κατά πόσον το πεδίο
\textit{τείνει} να συμπεριφέρεται τοπικά ως ακτινική πηγή ή καταβόθρα:
θετική απόκλιση σημαίνει ότι η γεωμετρική απόκλιση των δυναμικών γραμμών
υπερισχύει της σύγκλισης, ενώ αρνητική απόκλιση υποδηλώνει το αντίθετο.
Αυτό δε σημαίνει ότι εάν η απόκλιση είναι θετική σε κάποιο σημείο
τότε εκεί βρίσκεται πράγματι κάποια ακτινική πηγή,
παρά μόνον ότι η συνολική τοπική εκροή υπερβαίνει την εισροή
επειδή οι εκρέουσες δυναμικές γραμμές καταλαμβάνουν περισσότερο χώρο.
Ακτινική πηγή ή καταβόθρα είναι η οριακή περίπτωση όπου τελικά
όλες οι δυναμικές γραμμές εξέρχονται ή εισέρχονται
της επιφάνειας του στοιχείου όγκου καθώς αυτό
συρρικνώνεται αυθαίρετα γύρω από το σημείο·
ωστόσο, η θετική απόκλιση από μόνη της δεν αρκεί για να διακρίνει
μια ακτινική πηγή από την γενική περίπτωση,
δηλαδή τη συνδυαστική επενέργεια γεωμετρικής έκτασης των εκροών
και αυξημένων εξωστρεφών ταχυτήτων.

Η συζήτηση περί απόκλισης προϋποθέτει την παρουσία ρημάννειας μετρικής,
εφόσον γίνεται κατ' αρχήν λόγος περί όγκου.
Αυτό συνεχίζει να ισχύει για την θεώρηση της απόκλισης
υπό το πρίσμα της γεωμετρικής σύγκλισης ή απόκλισης των δυναμικών γραμμών
καθώς και της αυξομείωσης των ταχυτήτων
κατά μήκος της ροής του πεδίου.\footnote{Εάν εγκαταλείψουμε
την γεωμετρική οπτική της συμπεριφοράς
των δυναμικών γραμμών και της αυξομείωσης των ταχυτήτων,
τότε --με όρους διαφορικών μορφών--
η συζήτηση περί απόκλισης προϋποθέτει μόνο την παρουσία κάποιας μορφής όγκου,
αδιάφορο εάν αυτή επάγεται βάσει ρημάννειας μετρικής ή όχι.
Απέχουμε από αυτό το αφαιρετικό βήμα, εφόσον ενδιαφερόμαστε
για τη γεωμετρική ερμηνεία της απόκλισης
και για τις εφαρμογές της στη φυσική των πεδίων.}
Πράγματι, χωρίς κάποια έννοια μέτρησης αποστάσεων,
παραμένει ασαφές τι σημαίνει για μια δέσμη δυναμικών γραμμών
να αποκλίνουν
ή, ισοδύναμα, με ποιον τρόπο μπορούμε να παρακολουθήσουμε
την αύξηση της απόστασης δύο σωματιδίων
καθώς αυτά παρασύρονται από τη ροή του πεδίου·
παρομοίως, ελλείψει ρημάννειας μετρικής,
δεν υπάρχει εγγενής τρόπος σύγκρισης διανυσμάτων που ανήκουν
σε διαφορετικούς εφαπτομένους χώρους και άρα
δε μπορεί να γίνει λόγος για αυξομείωση της ταχύτητας
κατά μήκος των δυναμικών γραμμών του πεδίου.

%Υπ’ αυτήν την οπτική, πεδία με ταυτοτικά μηδενική απόκλιση
%επιδέχονται της ακόλουθης ερμηνείας:
%για κάθε δυναμική γραμμή,
%εάν θεωρήσουμε σωλήνα επαρκώς μικρής ακτίνας με κέντρο αυτήν,
%τότε η δέσμη των δυναμικών γραμμών που περιέχεται εντός του
%διατηρεί σταθερή εγκάρσια διατομή,
%εφόσον δεν υφίσταται συστολή ή διαστολή της ροής του πεδίου·
%τα σωματίδια που παρασύρονται από το πεδίο παραμένουν συνεπώς
%εγκλωβισμένα στον νοητό σωλήνα,
%πράγμα που δικαιολογεί την ονομασία των πεδίων
%με ταυτοτικά μηδενική απόκλιση ως \textit{σωληνοειδή}.

\section{Ρυθμός ογκομετρικής μεταβολής και συνολική εκροή πεδίου}

η μεταβολή του όγκου προέρχεται αποκλειστικά από την κανονική συνιστώσα της ταχύτητας πάνω στο σύνορο.

\begin{equation}
\left.\frac{d}{dt}\right|_{t=0}
\hspace{-10pt}
\operatorname{vol} \bigl(\phi_t(V)\bigr)
=
\int_{\partial V} \xi \cdot \hat{n} \,\, d\sigma
\end{equation}

%\section{Συνολική και κανονικοποιμένη ροή πεδίου}
%
%TODO

\section{Τυπικός ορισμός της απόκλισης}

TODO

\begin{equation}
\operatorname{div}\xi(p)
=
\lim_{\substack{V \to \{p\}}}
\frac{1}{\operatorname{vol} (V)}
\hspace{-0pt}
\int_{\partial V} \xi \cdot \hat{n} \,\, d\sigma
\end{equation}

\section{Θεώρημα Gauss-Ostrogradsky}

\begin{equation}
\int_V \operatorname{div}\xi(p) \,\, dv
=
\int_{\partial V} \xi \cdot \hat{n} \,\, d\sigma
\end{equation}

\section{Εξίσωση της συνέχειας}

TODO

\section{Έκφρασεις σε τοπικό σύστημα συντεταγμένων}

TODO

\begin{equation}
\nabla \cdot \mathbf{F}(\mathbf{r}) = 
\lim_{\Delta V \to 0} 
\frac{1}{\Delta V}
\oint_{\partial(\Delta V)} \mathbf{F} \cdot \hat{\mathbf{n}} \, dS
\end{equation}

\begin{equation}
\text{div} \, \boldsymbol{\xi}(\mathbf{p}) = 
\lim_{\Delta V \to 0} 
\frac{1}{\Delta V}
\oint_{\partial(\Delta V)} \boldsymbol{\xi} \cdot \hat{\mathbf{n}} \, d\sigma
\end{equation}

\begin{equation}
\text{div} \, \xi(p) = 
\lim_{\Delta V \to 0} 
\frac{1}{\Delta V}
\oint_{\partial(\Delta V)} \xi \cdot \hat{n} \, d\sigma
\end{equation}

\begin{equation}
\text{div} \, \xi(p) = 
\lim_{|V| \to 0} 
\frac{1}{|V|}
\oint_{\partial V} \xi^*(\hat{n}) \, d\sigma
\end{equation}


\chapter{Στροβιλισμός}

Ως \textit{στροβιλισμός} ενός διανυσματικού πεδίου εννοείται
το διανυσματικό μέγεθος που εκφράζει τοπική τάση του πεδίου
να προκαλεί ιδιοπεριστροφές,
δηλαδή να περιστρέφει τα στοιχειώδη χωρία του μέσου
γύρω από κάποιον στιγμαίο άξονά τους.
Πιο συγκεκριμένα, σε κάθε σημείο το χώρου,
ο στροβιλισμός ισούται με την σχετική ιδιοπεριστροφή υπό την επίδραση του πεδίου,
δηλαδή την στιγμιαία ιδιοπεριστροφή ενός ρέοντος χωρίου ανα μονάδα όγκου
καθώς αυτό συρρικνώνεται αυθαίρετα υπό το μελέτη σημείο.

\begin{equation}
\lim_{\substack{V \to \{p\}}}
\frac{1}{\operatorname{vol} (V)}
\hspace{-0pt}
\int_{\partial V} \hat{n} \times \xi \,\, d\sigma
\end{equation}


ο στροβιλισμός είναι ανάλογος της στερεοσωματικής περιστροφής
ενός στοιχείου όγκου υπό την επίδραση της ροής του πεδίου
καθώς αυτή συρρικνώνεται αυθαίρετα γύρω από το μελέτη σημείο.


\section{Ρυθμός κυκλοφοριακής ροής και μέσος εγκάρσιος όρος}
TODO

\section{Τυπικός ορισμός του στροβιλισμού}

TODO

\section{Θεώρημα Kelvin-Stokes}

TODO

\section{Εκφράσεις σε τοπικό σύστημα συντεταγμένων}

\chapter{Τοπική Προσέγγιση Εμβαδού και Σύγκλιση Επιφανειακών Ολοκληρωμάτων}

\section{Ασυμπτωτικός Συμβολισμός}

Πριν προχωρήσουμε στην κυρίως ανάλυση, υπενθυμίζουμε τον ακριβή ορισμό του συμβολισμού \(O\) (μεγάλο όμικρον) που θα χρησιμοποιείται σε όλο το Παράρτημα.

\begin{definition}[Συμβολισμός \(O\) για μικρές μεταβολές]
Έστω \(E(\Delta u, \Delta v)\) πραγματική συνάρτηση ορισμένη σε μια περιοχή του \((0,0)\) και \(p \ge 0\) ακέραιος. Γράφουμε
\[
E(\Delta u, \Delta v) = O(h^p), \quad h = \max(|\Delta u|, |\Delta v|),
\]
αν υπάρχουν σταθερές \(C > 0\) και \(\delta > 0\) ώστε για όλα τα \(\Delta u, \Delta v\) με \(0 < h < \delta\) να ισχύει
\[
|E(\Delta u, \Delta v)| \le C \cdot h^p.
\]
\end{definition}

\begin{remark}
Στην περίπτωση του υπολοίπου Taylor, όπου η μεταβλητή απόστασης είναι \(\rho = \sqrt{\delta u^2 + \delta v^2}\), χρησιμοποιούμε τη σχέση \(\rho \le \sqrt{2} h\) για να μεταφράσουμε το φράγμα από \(O(\rho^2)\) σε \(O(h^2)\). Πράγματι, αν \(|R_1| \le M \rho^2\), τότε \(|R_1| \le M \cdot 2h^2 = (2M) h^2\), άρα \(R_1 = O(h^2)\).

Σε καθολικές εκτιμήσεις, όπου έχουμε διαμέριση του χωρίου \(D\) σε \(N\) πλακίδια, το \(h\) δηλώνει τη μέγιστη διάμετρο των πλακιδίων: \(h = \max_i \operatorname{diam}(Q_i)\). Το πλήθος των πλακιδίων \(N\) ικανοποιεί \(N = O(h^{-2})\), δηλαδή υπάρχει σταθερά \(C'\) ώστε \(N \le C' h^{-2}\).
\end{remark}

\section{Γεωμετρική Εικόνα και Βασικοί Ορισμοί}

Έστω μια κανονική \(C^2\)-παραμέτρηση επιφάνειας
\[
\mathbf{r} : D \subset \mathbb{R}^2 \to \mathbb{R}^3, \quad (u,v) \mapsto \mathbf{r}(u,v),
\]
όπου \(D\) είναι ανοικτό και φραγμένο χωρίο. Η κανονικότητα σημαίνει
\[
\mathbf{r}_u \times \mathbf{r}_v \neq \mathbf{0}, \quad \forall (u,v) \in D.
\]

Θεωρούμε σημείο \(p_0 = \mathbf{r}(u_0, v_0) \in \mathcal{S}\) και θετικές μεταβολές \(\Delta u, \Delta v\) αρκετά μικρές ώστε το ορθογώνιο
\[
R = [u_0, u_0+\Delta u] \times [v_0, v_0+\Delta v]
\]
να περιέχεται εξ ολοκλήρου στο \(D\).

\begin{definition}[Στοιχειώδη Πλακίδια]
\begin{enumerate}
\item \textbf{Παραμετρικό πλακίδιο} \(Q\): η εικόνα του ορθογωνίου \(R\) μέσω της \(\mathbf{r}\),
\[
Q = \mathbf{r}(R) = \{ \mathbf{r}(u,v) : (u,v) \in R \}.
\]
\item \textbf{Εφαπτόμενο πλακίδιο} \(P\): το παραλληλόγραμμο στον εφαπτόμενο χώρο \(T_{p_0}\mathcal{S}\) που παράγεται από τα διανύσματα
\[
\mathbf{r}_u(p_0)\Delta u, \quad \mathbf{r}_v(p_0)\Delta v.
\]
\end{enumerate}
\end{definition}

\begin{definition}[Συναρτήσεις Μέτρου]
\begin{enumerate}
\item \textbf{Συνάρτηση εμβαδικού στοιχείου:}
\[
f(u,v) = \| \mathbf{r}_u(u,v) \times \mathbf{r}_v(u,v) \|.
\]
\item \textbf{Γενικευμένη συνάρτηση μέτρου για βαθμωτό πεδίο} \(g: \mathcal{S} \to \mathbb{R}\):
\[
F(u,v) = g(\mathbf{r}(u,v)) \cdot f(u,v).
\]
\end{enumerate}
\end{definition}

Εφόσον η \(\mathbf{r}\) είναι \(C^2\) και \(f > 0\) (από κανονικότητα), η \(f\) είναι \(C^2\) στο \(D\). Αν επιπλέον η \(g \circ \mathbf{r}\) είναι \(C^2\), τότε και η \(F\) είναι \(C^2\).

\section{Τοπική Προσέγγιση του Εμβαδού}

Το πραγματικό εμβαδόν του καμπυλόγραμμου πλακιδίου \(Q\) είναι
\[
A(Q) = \iint_Q d\mathcal{A} = \int_{v_0}^{v_0+\Delta v} \int_{u_0}^{u_0+\Delta u} f(u,v) \, du \, dv.
\]

Το εμβαδόν του εφαπτόμενου πλακιδίου \(P\) είναι
\[
A(P) = \| \mathbf{r}_u \times \mathbf{r}_v \|_{p_0} \, \Delta u \, \Delta v = f(p_0) \, \Delta u \, \Delta v.
\]

\begin{theorem}[Τοπικό σφάλμα για το εμβαδόν]
\label{thm:local_area}
Θέτουμε \(h = \max(\Delta u, \Delta v)\). Τότε
\[
A(Q) = A(P) + O(h^3).
\]
\end{theorem}

\begin{proof}
Αναπτύσσουμε την \(f\) κατά Taylor γύρω από το \(p_0 = (u_0, v_0)\). Για \(\delta u = u - u_0, \delta v = v - v_0\),
\[
f(u,v) = f(p_0) + f_u(p_0)\,\delta u + f_v(p_0)\,\delta v + R_1(\delta u, \delta v),
\]
όπου το υπόλοιπο \(R_1\) ορίζεται ως η διαφορά της \(f\) από το γραμμικό της μέρος:
\[
R_1(\delta u, \delta v) = f(u,v) - \left[ f(p_0) + f_u(p_0)\delta u + f_v(p_0)\delta v \right].
\]

Εφόσον η \(f\) είναι \(C^2\) στο \(D\), οι δεύτερες παράγωγοί της είναι συνεχείς και φραγμένες σε κάθε συμπαγές υποσύνολο του \(D\). Από το θεώρημα Taylor για \(C^2\) συναρτήσεις, υπάρχει σταθερά \(M\) (ανεξάρτητη των \(\delta u, \delta v\) για μικρές μεταβολές) ώστε
\[
|R_1(\delta u, \delta v)| \le M \cdot (\delta u^2 + \delta v^2).
\]
Ισοδύναμα, με \(\rho = \sqrt{\delta u^2 + \delta v^2}\),
\[
R_1(\delta u, \delta v) = O(\rho^2).
\]

Μετατοπίζουμε το χωρίο ολοκλήρωσης στην αρχή και ολοκληρώνουμε:
\[
A(Q) = \int_0^{\Delta v} \int_0^{\Delta u} \left[ f(p_0) + f_u(p_0)\,\delta u + f_v(p_0)\,\delta v + R_1 \right] d\delta u \, d\delta v.
\]

Υπολογίζουμε κάθε όρο χωριστά.

\begin{enumerate}
\item \textbf{Σταθερός όρος:}
\[
\int_0^{\Delta v} \int_0^{\Delta u} f(p_0) \, d\delta u \, d\delta v = f(p_0) \Delta u \Delta v = A(P).
\]

\item \textbf{Γραμμικοί όροι:}
\[
f_u(p_0) \int_0^{\Delta v} \int_0^{\Delta u} \delta u \, d\delta u \, d\delta v = f_u(p_0) \cdot \Delta v \cdot \frac{(\Delta u)^2}{2},
\]
\[
f_v(p_0) \int_0^{\Delta v} \int_0^{\Delta u} \delta v \, d\delta u \, d\delta v = f_v(p_0) \cdot \Delta u \cdot \frac{(\Delta v)^2}{2}.
\]
Αυτοί είναι τάξης \(O(h^3)\), διότι π.χ. \(|\Delta v| (\Delta u)^2 \le h \cdot h^2 = h^3\).

\item \textbf{Υπόλοιπο:}
Λόγω της \(R_1 = O(\rho^2)\) και της \(\rho \le \sqrt{2}h\) στο χωρίο ολοκλήρωσης, έχουμε \(R_1 = O(h^2)\) (βλ. Παρατήρηση Α.1). Το διπλό ολοκλήρωμα του υπολοίπου είναι
\[
\int_0^{\Delta v} \int_0^{\Delta u} O(h^2) \, d\delta u \, d\delta v = O(h^2) \cdot \Delta u \Delta v = O(h^2) \cdot O(h^2) = O(h^4),
\]
αμελητέο συγκρινόμενο με το \(O(h^3)\).
\end{enumerate}

Επομένως,
\[
A(Q) - A(P) = f_u(p_0) \frac{(\Delta u)^2 \Delta v}{2} + f_v(p_0) \frac{\Delta u (\Delta v)^2}{2} + O(h^4) = O(h^3).
\]
\end{proof}

\textbf{Γεωμετρική ερμηνεία.} Το κύριο μέρος του σφάλματος προέρχεται από τη μεταβολή του μήκους του καθέτου διανύσματος κατά μήκος του πλακιδίου (πρώτες παράγωγοι της \(f\)). Αυτή η μεταβολή ολοκληρώνεται και δίνει όρους τάξης \(h^3\), δηλαδή μία τάξη μικρότερους από το ίδιο το εμβαδόν (\(h^2\)). Οι δεύτερες παράγωγοι (καμπυλότητα) συνεισφέρουν μόνο σε όρους \(h^4\) και άνω.

\section{Γενίκευση: Επιφανειακό Ολοκλήρωμα Βαθμωτής Συνάρτησης}

\subsection{Τοπική Προσέγγιση}

Έστω \(g: \mathcal{S} \to \mathbb{R}\) βαθμωτή συνάρτηση τέτοια ώστε η σύνθεση \(g \circ \mathbf{r}\) να είναι \(C^2\) στο \(D\). Το επιφανειακό ολοκλήρωμα της \(g\) στο πλακίδιο \(Q\) είναι
\[
I_g(Q) = \iint_Q g \, d\mathcal{A} = \int_{v_0}^{v_0+\Delta v} \int_{u_0}^{u_0+\Delta u} g(\mathbf{r}(u,v)) \, f(u,v) \, du \, dv.
\]

Η διακριτή προσέγγιση χρησιμοποιεί την τιμή της \(g\) στο σημείο βάσης \(p_0\) και το εμβαδόν του εφαπτόμενου πλακιδίου:
\[
I_g(P) = g(p_0) \cdot A(P) = g(p_0) \, f(p_0) \, \Delta u \, \Delta v.
\]

\begin{theorem}[Τοπικό σφάλμα για βαθμωτή συνάρτηση]
\label{thm:local_scalar}
Έστω \(F(u,v) = g(\mathbf{r}(u,v)) \cdot f(u,v)\) και υποθέτουμε ότι η \(F\) είναι \(C^2\) στο \(D\). Τότε, με \(h = \max(\Delta u, \Delta v)\),
\[
I_g(Q) = I_g(P) + O(h^3).
\]
\end{theorem}

\begin{proof}
Η απόδειξη είναι πανομοιότυπη με αυτήν του Θεωρήματος \ref{thm:local_area}, αντικαθιστώντας την \(f\) με την \(F\). Αναπτύσσουμε την \(F\) κατά Taylor γύρω από το \(p_0\):
\[
F(u,v) = F(p_0) + F_u(p_0)\,\delta u + F_v(p_0)\,\delta v + R_1^F(\delta u, \delta v),
\]
όπου \(F(p_0) = g(p_0) f(p_0)\). Λόγω της \(C^2\)-ομαλότητας, \(R_1^F = O(\rho^2)\). Ολοκληρώνοντας, ο σταθερός όρος δίνει \(I_g(P)\), οι γραμμικοί όροι δίνουν \(O(h^3)\) και το υπόλοιπο \(O(h^4)\).
\end{proof}

\subsection{Καθολική Σύγκλιση}

Θεωρούμε διαμέριση του \(D\) σε \(N\) ορθογώνια \(Q_i\) (\(i=1,\dots,N\)), με πλευρές \(\Delta u_i, \Delta v_i\). Η μέγιστη διάμετρος της διαμέρισης είναι
\[
h = \max_i \operatorname{diam}(Q_i) = \max_i \max(\Delta u_i, \Delta v_i).
\]

Σε κάθε \(Q_i\) επιλέγουμε την κάτω αριστερή κορυφή \((u_i, v_i)\) και ορίζουμε το διακριτό άθροισμα
\[
I_h(g) = \sum_{i=1}^N g(\mathbf{r}(u_i, v_i)) \, f(u_i, v_i) \, \Delta u_i \Delta v_i.
\]

\begin{theorem}[Σύγκλιση επιφανειακού ολοκληρώματος βαθμωτής συνάρτησης]
\label{thm:global_scalar}
Υπό την προϋπόθεση ότι η \(F = (g \circ \mathbf{r}) \cdot f\) είναι \(C^2\) στο \(D\),
\[
\lim_{h \to 0} I_h(g) = \iint_{\mathcal{S}} g \, d\mathcal{A} = \iint_D g(\mathbf{r}(u,v)) \, f(u,v) \, du \, dv.
\]
\end{theorem}

\begin{proof}
Το συνολικό σφάλμα είναι
\[
\mathcal{E}_h = \left| \iint_{\mathcal{S}} g \, d\mathcal{A} - I_h(g) \right| \le \sum_{i=1}^N \left| I_g(Q_i) - I_g(P_i) \right|.
\]

Από το Θεώρημα \ref{thm:local_scalar}, για κάθε πλακίδιο υπάρχει σταθερά \(K\) (ανεξάρτητη του \(i\), λόγω της \(C^2\)-ομαλότητας της \(F\) στο συμπαγές \(\overline{D}\)) ώστε
\[
|E_i| = |I_g(Q_i) - I_g(P_i)| \le K (\operatorname{diam} Q_i)^3 \le K h^3.
\]

Το πλήθος \(N\) των πλακιδίων ικανοποιεί \(N = O(h^{-2})\), διότι το συνολικό εμβαδόν του \(D\) είναι πεπερασμένο και κάθε \(Q_i\) έχει εμβαδόν \(\Delta u_i \Delta v_i = O(h^2)\). Αυστηρότερα, υπάρχει σταθερά \(C\) ώστε \(N \le C h^{-2}\) για όλες τις αρκετά λεπτές διαμερίσεις.

Συνεπώς,
\[
\mathcal{E}_h \le \sum_{i=1}^N K h^3 = N \cdot K h^3 \le (C h^{-2}) \cdot K h^3 = CK h.
\]

Καθώς \(h \to 0\), το δεξιό μέλος τείνει στο μηδέν.
\end{proof}

\begin{corollary}[Σύγκλιση για το εμβαδόν]
Για \(g \equiv 1\), έχουμε \(F = f\) και το διακριτό άθροισμα
\[
A_h(\mathcal{S}) = \sum_{i=1}^N f(u_i, v_i) \Delta u_i \Delta v_i
\]
συγκλίνει στο εμβαδόν
\[
A(\mathcal{S}) = \iint_D f(u,v) \, du \, dv.
\]
\end{corollary}

\subsection{Ανεξαρτησία από την Παραμέτρηση}

Το επιφανειακό ολοκλήρωμα που ορίσαμε είναι ανεξάρτητο από την επιλογή της \(C^2\)-παραμέτρησης. Πράγματι, αν \(\tilde{\mathbf{r}} : \tilde{D} \to \mathcal{S}\) είναι μια άλλη κανονική \(C^2\)-παραμέτρηση, η απεικόνιση αλλαγής συντεταγμένων \(\phi = \tilde{\mathbf{r}}^{-1} \circ \mathbf{r} : D \to \tilde{D}\) είναι \(C^1\)-αμφιδιαφόριση. Το θεώρημα αλλαγής μεταβλητών για πολλαπλά ολοκληρώματα δίνει
\[
\iint_D g(\mathbf{r}(u,v)) \|\mathbf{r}_u \times \mathbf{r}_v\| \, du \, dv = \iint_{\tilde{D}} g(\tilde{\mathbf{r}}(\tilde{u},\tilde{v})) \|\tilde{\mathbf{r}}_{\tilde{u}} \times \tilde{\mathbf{r}}_{\tilde{v}}\| \, d\tilde{u} \, d\tilde{v},
\]
αφού η Ιακωβιανή ορίζουσα της αλλαγής συντεταγμένων ακριβώς αντισταθμίζει τον μετασχηματισμό του στοιχείου εμβαδού. Επομένως το ολοκλήρωμα είναι γεωμετρική ποσότητα, ανεξάρτητη της παραμετρικοποίησης.

\section{Επιφανειακό Ολοκλήρωμα Διανυσματικής Συνάρτησης}

\subsection{Γενικός Ορισμός}

Για μια διανυσματική συνάρτηση \(\mathbf{F} : \mathcal{S} \to \mathbb{R}^3\), το επιφανειακό ολοκλήρωμα ορίζεται κατά συνιστώσες:
\[
\iint_{\mathcal{S}} \mathbf{F} \, d\mathcal{A} = \left( \iint_{\mathcal{S}} F_x \, d\mathcal{A}, \; \iint_{\mathcal{S}} F_y \, d\mathcal{A}, \; \iint_{\mathcal{S}} F_z \, d\mathcal{A} \right).
\]

Κάθε συνιστώσα είναι βαθμωτό επιφανειακό ολοκλήρωμα και επομένως όλη η ανάλυση της §Α.4 εφαρμόζεται αυτούσια. Το διακριτό ανάλογο είναι το διανυσματικό άθροισμα
\[
\mathbf{I}_h(\mathbf{F}) = \sum_{i=1}^N \mathbf{F}(\mathbf{r}(u_i, v_i)) \, f(u_i, v_i) \, \Delta u_i \Delta v_i,
\]
και συγκλίνει συνιστωσικά στο ολοκλήρωμα καθώς \(h \to 0\), υπό την προϋπόθεση ότι κάθε συνιστώσα της \(\mathbf{F} \circ \mathbf{r}\) είναι \(C^2\) ώστε οι αντίστοιχες \(F\) συναρτήσεις να είναι \(C^2\).

\subsection{Το Ολοκλήρωμα του Μοναδιαίου Καθέτου Διανύσματος}

Μια θεμελιώδης ειδική περίπτωση προκύπτει όταν \(\mathbf{F} = \mathbf{n}\), το μοναδιαίο κάθετο διάνυσμα της επιφάνειας. Αυτό ορίζεται από τη σχέση
\[
\mathbf{n}(u,v) = \frac{\mathbf{r}_u \times \mathbf{r}_v}{\|\mathbf{r}_u \times \mathbf{r}_v\|} = \frac{\mathbf{r}_u \times \mathbf{r}_v}{f(u,v)}.
\]

Το επιφανειακό ολοκλήρωμά του απλοποιείται δραματικά:
\[
\iint_{\mathcal{S}} \mathbf{n} \, d\mathcal{A} = \iint_D \mathbf{n}(\mathbf{r}(u,v)) \, f(u,v) \, du \, dv = \iint_D (\mathbf{r}_u \times \mathbf{r}_v) \, du \, dv.
\]

Παρατηρούμε ότι το βαθμωτό μέτρο \(f\) απαλείφεται εντελώς και το ολοκλήρωμα ανάγεται στο ολοκλήρωμα του μη μοναδιαίου καθέτου διανύσματος \(\mathbf{N} = \mathbf{r}_u \times \mathbf{r}_v\).

Η διακριτή προσέγγιση σε κάθε πλακίδιο είναι
\[
\mathbf{n}_i \cdot A(P_i) = \frac{\mathbf{r}_u \times \mathbf{r}_v}{f} \cdot f \, \Delta u_i \Delta v_i = (\mathbf{r}_u \times \mathbf{r}_v) \, \Delta u_i \Delta v_i,
\]
δηλαδή ακριβώς το διάνυσμα του εφαπτόμενου παραλληλογράμμου (κάθετο με μέτρο ίσο με το εμβαδόν του). Το συνολικό διακριτό άθροισμα
\[
\mathbf{K}_h = \sum_{i=1}^N (\mathbf{r}_u \times \mathbf{r}_v)_{(u_i, v_i)} \, \Delta u_i \Delta v_i
\]
συγκλίνει, σύμφωνα με τη γενική θεωρία, στο
\[
\mathbf{K} = \iint_D (\mathbf{r}_u \times \mathbf{r}_v) \, du \, dv.
\]

\subsection{Γεωμετρική Ερμηνεία}

Το διάνυσμα
\[
\mathbf{K} = \iint_{\mathcal{S}} \mathbf{n} \, d\mathcal{A}
\]
έχει πλούσιο γεωμετρικό περιεχόμενο.

\begin{enumerate}
\item \textbf{Επίπεδο χωρίο.} Αν η \(\mathcal{S}\) περιέχεται σε ένα επίπεδο με σταθερό μοναδιαίο κάθετο \(\mathbf{n}_0\), τότε
\[
\mathbf{K} = \mathbf{n}_0 \iint_{\mathcal{S}} d\mathcal{A} = \mathbf{n}_0 \cdot A(\mathcal{S}).
\]
Πρόκειται για ένα διάνυσμα κάθετο στο επίπεδο, με μέτρο ίσο με το εμβαδόν του χωρίου.

\item \textbf{Κλειστή επιφάνεια.} Αν η \(\mathcal{S}\) είναι κλειστή (συμπαγής και χωρίς σύνορο) και περικλείει όγκο \(V\), τότε
\[
\mathbf{K} = \iint_{\mathcal{S}} \mathbf{n} \, d\mathcal{A} = \mathbf{0}.
\]
Η απόδειξη είναι μια άμεση εφαρμογή του θεωρήματος της απόκλισης: για οποιοδήποτε σταθερό διάνυσμα \(\mathbf{c}\),
\[
\mathbf{c} \cdot \iint_{\mathcal{S}} \mathbf{n} \, d\mathcal{A} = \iint_{\mathcal{S}} \mathbf{c} \cdot \mathbf{n} \, d\mathcal{A} = \iiint_V (\nabla \cdot \mathbf{c}) \, dV = \iiint_V 0 \, dV = 0.
\]
Εφόσον αυτό ισχύει για κάθε \(\mathbf{c}\), το διάνυσμα \(\mathbf{K}\) είναι υποχρεωτικά μηδενικό.

\item \textbf{Ανοικτή επιφάνεια με σύνορο.} Όταν η \(\mathcal{S}\) είναι ανοικτή, το διάνυσμα \(\mathbf{K}\) δεν είναι κατ' ανάγκη μηδέν. Οι συνιστώσες του έχουν την εξής ερμηνεία:
\[
K_x = \iint_{\mathcal{S}} n_x \, d\mathcal{A}, \quad K_y = \iint_{\mathcal{S}} n_y \, d\mathcal{A}, \quad K_z = \iint_{\mathcal{S}} n_z \, d\mathcal{A}.
\]
Από τον ορισμό του μοναδιαίου καθέτου, η συνιστώσα \(n_x\) είναι το συνημίτονο της γωνίας που σχηματίζει το κάθετο διάνυσμα με τον άξονα \(x\). Το ολοκλήρωμα \(\iint_{\mathcal{S}} n_x \, d\mathcal{A}\) ισούται ακριβώς με το \textbf{προσημασμένο εμβαδόν της προβολής} της \(\mathcal{S}\) στο επίπεδο \(yz\). Το πρόσημο είναι θετικό αν το κάθετο διάνυσμα ``βλέπει'' προς τη θετική κατεύθυνση \(x\), και αρνητικό διαφορετικά. Ανάλογα για τις άλλες συνιστώσες.

Συνεπώς, το \(\mathbf{K}\) είναι ένα διάνυσμα του οποίου οι συνιστώσες είναι τα προσημασμένα εμβαδά των προβολών της επιφάνειας στα επίπεδα συντεταγμένων.
\end{enumerate}

\subsection{Σύνδεση με το Θεώρημα του Stokes}

Η παραπάνω ερμηνεία συνδέεται άμεσα με το θεώρημα του Stokes (στη κλασική διανυσματική του μορφή, θεώρημα Kelvin--Stokes). Για ένα διανυσματικό πεδίο \(\mathbf{F}\),
\[
\iint_{\mathcal{S}} (\nabla \times \mathbf{F}) \cdot \mathbf{n} \, d\mathcal{A} = \oint_{\partial \mathcal{S}} \mathbf{F} \cdot d\mathbf{r}.
\]

Επιλέγοντας \(\mathbf{F}\) τέτοιο ώστε \(\nabla \times \mathbf{F} = \mathbf{c}\) (ένα σταθερό διάνυσμα), το αριστερό μέλος γίνεται \(\mathbf{c} \cdot \iint_{\mathcal{S}} \mathbf{n} \, d\mathcal{A}\), ενώ το δεξιό είναι ένα επικαμπύλιο ολοκλήρωμα που εξαρτάται μόνο από το σύνορο \(\partial\mathcal{S}\). Αυτό αποδεικνύει ότι το \(\iint_{\mathcal{S}} \mathbf{n} \, d\mathcal{A}\) είναι μια ποσότητα που εξαρτάται αποκλειστικά από το σύνορο της \(\mathcal{S}\) --- γεγονός που εξηγεί γιατί μηδενίζεται για κλειστές επιφάνειες (όπου \(\partial\mathcal{S} = \emptyset\)) και γιατί για ανοικτές επιφάνειες σχετίζεται με τις προβολικές ιδιότητες του συνόρου.

\section{Ο Θεμελιώδης Μηχανισμός Σύγκλισης}

Όλες οι παραπάνω αποδείξεις σύγκλισης βασίζονται σε έναν ενιαίο, θεμελιώδη μηχανισμό που αξίζει να αναδειχθεί ρητά.

\begin{theorem}[Γενική Αρχή Σύγκλισης]
\label{thm:general_principle}
Έστω \(H(u,v)\) μια \(C^2\) συνάρτηση στο φραγμένο χωρίο \(D\) και έστω διαμέριση του \(D\) σε \(N\) ορθογώνια \(Q_i\) μέγιστης διαμέτρου \(h\). Τότε
\[
\left| \iint_D H(u,v) \, du \, dv - \sum_{i=1}^N H(u_i, v_i) \Delta u_i \Delta v_i \right| = O(h),
\]
όπου \((u_i, v_i)\) είναι η κάτω αριστερή κορυφή του \(Q_i\).
\end{theorem}

\begin{proof}
Το σφάλμα σε κάθε πλακίδιο είναι
\[
E_i = \iint_{Q_i} H(u,v) \, du \, dv - H(u_i, v_i) \Delta u_i \Delta v_i.
\]

Αναπτύσσοντας την \(H\) γύρω από το \((u_i, v_i)\) ακριβώς όπως στα Θεωρήματα \ref{thm:local_area} και \ref{thm:local_scalar}, βρίσκουμε \(E_i = O(h^3)\) (οι γραμμικοί όροι δίνουν \(O(h^3)\) και το υπόλοιπο \(O(h^4)\)). Το πλήθος των πλακιδίων είναι \(N = O(h^{-2})\). Άρα το συνολικό σφάλμα είναι
\[
\sum_{i=1}^N |E_i| \le N \cdot O(h^3) = O(h^{-2}) \cdot O(h^3) = O(h) \to 0.
\]
\end{proof}

Αυτή η αρχή είναι ο λόγος που η μέθοδος των αθροισμάτων Riemann λειτουργεί για επιφανειακά ολοκληρώματα: το τοπικό σφάλμα γραμμικής προσέγγισης είναι μία τάξη μικρότερο (\(h^3\)) από το στοιχειώδες εμβαδόν (\(h^2\)), με αποτέλεσμα η συσσώρευσή του σε \(O(h^{-2})\) πλακίδια να δίνει συνολικό σφάλμα \(O(h)\), το οποίο εξαφανίζεται στο όριο \(h \to 0\).

\section{Συμπέρασμα}

Σε αυτό το Παράρτημα αποδείξαμε με αυστηρό τρόπο τη σύγκλιση της διακριτής προσέγγισης των επιφανειακών ολοκληρωμάτων μέσω εφαπτόμενων παραλληλογράμμων. Τα κύρια σημεία είναι:

\begin{enumerate}
\item Το τοπικό σφάλμα μεταξύ του πραγματικού εμβαδού ενός μικρού καμπυλόγραμμου πλακιδίου και του εφαπτόμενου παραλληλογράμμου του είναι \(O(h^3)\), όπου \(h\) είναι η μέγιστη διάσταση του πλακιδίου.

\item Το ίδιο ισχύει για το ολοκλήρωμα οποιασδήποτε βαθμωτής ή διανυσματικής συνάρτησης, αρκεί η συνάρτηση προς ολοκλήρωση (πολλαπλασιασμένη με το στοιχείο εμβαδού) να είναι αρκετά ομαλή (\(C^2\)).

\item Το συνολικό σφάλμα μιας διαμέρισης με \(N = O(h^{-2})\) πλακίδια είναι \(O(h)\), το οποίο μηδενίζεται στο όριο \(h \to 0\), εξασφαλίζοντας τη σύγκλιση του διακριτού αθροίσματος στο συνεχές ολοκλήρωμα.

\item Το ολοκλήρωμα του μοναδιαίου καθέτου διανύσματος \(\iint_{\mathcal{S}} \mathbf{n} \, d\mathcal{A}\) έχει ιδιαίτερη γεωμετρική σημασία: οι συνιστώσες του είναι τα προσημασμένα εμβαδά των προβολών της επιφάνειας, μηδενίζεται για κλειστές επιφάνειες, και συνδέεται με το θεώρημα του Stokes.
\end{enumerate}

Αυτά τα αποτελέσματα θεμελιώνουν τον κλασικό ορισμό του επιφανειακού ολοκληρώματος στη διαφορική γεωμετρία και τη φυσική, και δικαιολογούν γιατί η γραμμική προσέγγιση του εφαπτόμενου επιπέδου επαρκεί για τον ακριβή υπολογισμό γεωμετρικών και φυσικών ποσοτήτων πάνω σε καμπύλες επιφάνειες.

\appendix
\chapter{Appendix A}

\section{Γεωμετρικές ερμηνείες εξωτερικού γινομένου}

\begin{equation}
\xi
=
\langle \hspace{1.5pt} \xi, \hat{n} \hspace{1.5pt} \rangle \hspace{1.5pt} \hat{n}
\hspace{1pt}
+
\hspace{1pt}
\hat{n} \times (\xi \times \hat{n})
\end{equation}

\noindent
Εγκάρσια συνιστώσα:

\begin{equation}
\xi_{tr} = \hat{n} \times (\xi \times \hat{n})
\end{equation}

\noindent
Iδιοπεριστροφή (``spin'') σημείου· κωδικοποιεί πληροφορία για τον άξονα
της στιγμαίας περιστροφικής κίνησης τπου επάγεται από την εγκάρσια
συνιστώσα (δηλαδή το επίπεδο στο οποίο λαμβάνει χώρα η στιγμαία περιστροφή)
καθώς και το μέτρο της αντίστοιχης γωνιακής ταχύτητας:

\begin{equation}
\operatorname{spin}(p, \xi) = \hat{n}(p) \times \xi(p)
\end{equation}

\noindent
Ιδιοπεριστροφή (``spin'') χωρίου: η συνισταμένη του spin όλων των σημείων του·
εφόσον τα εσωτερικά σημεία δεν βρίσκονται στην επιφάνεια και άρα δεν υφίστανται
απομονωμένη την περιστροφική επίδραση του πεδίου:

\begin{equation}
\int_{\partial V} \hspace{-0.5pt} \hat{n} \times \xi \,\, d\sigma
\end{equation}


\section{Taylor}

\begin{lemma}[Φράγμα Υπολοίπου Taylor Πρώτης Τάξης για \(C^2\) Συναρτήσεις]
\label{lem:C2_Taylor_bound_Rn}
Έστω \(D \subset \mathbb{R}^n\) ανοικτό χωρίο και \(f \in C^2(D)\) (δηλαδή η \(f\) έχει συνεχείς δεύτερες μερικές παραγώγους στο \(D\)). 

Για κάθε συμπαγές υποσύνολο \(K \subset D\), υπάρχει σταθερά \(M_K > 0\) (που εξαρτάται μόνο από την \(f\) και το \(K\)) με την ακόλουθη ιδιότητα:

Για κάθε σημείο \(\mathbf{x}_0 = (x_1^0, \dots, x_n^0) \in K\) και για κάθε \(\mathbf{h} = (h_1, \dots, h_n) \in \mathbb{R}^n\) τέτοιο ώστε το ευθύγραμμο τμήμα από το \(\mathbf{x}_0\) μέχρι το \(\mathbf{x}_0 + \mathbf{h}\) να περιέχεται στο \(K\), το υπόλοιπο πρώτης τάξης
\[
R_1(\mathbf{h}) = f(\mathbf{x}_0 + \mathbf{h}) - \left[ f(\mathbf{x}_0) + \sum_{i=1}^n \frac{\partial f}{\partial x_i}(\mathbf{x}_0) \, h_i \right]
\]
ικανοποιεί το φράγμα
\[
|R_1(\mathbf{h})| \le M_K \cdot \|\mathbf{h}\|^2,
\]
όπου \(\|\mathbf{h}\| = \sqrt{h_1^2 + \cdots + h_n^2}\) είναι η Ευκλείδεια νόρμα.

Ισοδύναμα, ισχύει
\[
R_1(\mathbf{h}) = O(\|\mathbf{h}\|^2) \quad \text{καθώς } \mathbf{h} \to \mathbf{0}.
\]
\end{lemma}

\begin{proof}
Η απόδειξη βασίζεται στο θεώρημα Taylor με ολοκληρωτικό υπόλοιπο για συναρτήσεις πολλών μεταβλητών.

Θεωρούμε τη βοηθητική συνάρτηση μιας μεταβλητής
\[
\varphi(t) = f(\mathbf{x}_0 + t\,\mathbf{h}), \quad t \in [0,1].
\]

Αφού η \(f\) είναι \(C^2\) στο \(D\) και το ευθύγραμμο τμήμα από το \(\mathbf{x}_0\) μέχρι το \(\mathbf{x}_0 + \mathbf{h}\) περιέχεται στο \(K \subset D\), η \(\varphi\) είναι καλά ορισμένη και \(C^2\) στο \([0,1]\).

Υπολογίζουμε τις παραγώγους της \(\varphi\) χρησιμοποιώντας τον κανόνα αλυσίδας.

\textbf{Πρώτη παράγωγος:}
\[
\varphi'(t) = \sum_{i=1}^n \frac{\partial f}{\partial x_i}(\mathbf{x}_0 + t\mathbf{h}) \cdot h_i = \nabla f(\mathbf{x}_0 + t\mathbf{h}) \cdot \mathbf{h}.
\]

\textbf{Δεύτερη παράγωγος:}
\begin{align*}
\varphi''(t) &= \frac{d}{dt} \left( \sum_{i=1}^n \frac{\partial f}{\partial x_i}(\mathbf{x}_0 + t\mathbf{h}) \, h_i \right) \\
&= \sum_{i=1}^n \left( \sum_{j=1}^n \frac{\partial^2 f}{\partial x_j \partial x_i}(\mathbf{x}_0 + t\mathbf{h}) \, h_j \right) h_i \\
&= \sum_{i=1}^n \sum_{j=1}^n \frac{\partial^2 f}{\partial x_i \partial x_j}(\mathbf{x}_0 + t\mathbf{h}) \, h_i \, h_j.
\end{align*}

Παρατηρούμε ότι η \(\varphi''(t)\) είναι μια τετραγωνική μορφή στον \(\mathbf{h}\) με πίνακα τον Εσσιανό \(H_f(\mathbf{x}_0 + t\mathbf{h}) = \left[ \frac{\partial^2 f}{\partial x_i \partial x_j} \right]_{i,j=1}^n\):
\[
\varphi''(t) = \mathbf{h}^T \, H_f(\mathbf{x}_0 + t\mathbf{h}) \, \mathbf{h}.
\]

Το ανάπτυγμα Taylor της \(\varphi\) γύρω από το \(t=0\) με ολοκληρωτικό υπόλοιπο πρώτης τάξης δίνει:
\[
\varphi(1) = \varphi(0) + \varphi'(0) + \int_0^1 (1-t) \, \varphi''(t) \, dt.
\]

Αντικαθιστούμε:
\begin{align*}
\varphi(0) &= f(\mathbf{x}_0), \\
\varphi'(0) &= \sum_{i=1}^n \frac{\partial f}{\partial x_i}(\mathbf{x}_0) \, h_i, \\
\varphi(1) &= f(\mathbf{x}_0 + \mathbf{h}).
\end{align*}

Συνεπώς:
\[
f(\mathbf{x}_0 + \mathbf{h}) = f(\mathbf{x}_0) + \sum_{i=1}^n \frac{\partial f}{\partial x_i}(\mathbf{x}_0) \, h_i + \int_0^1 (1-t) \, \varphi''(t) \, dt.
\]

Άρα το υπόλοιπο πρώτης τάξης ταυτίζεται με το ολοκλήρωμα:
\[
R_1(\mathbf{h}) = \int_0^1 (1-t) \, \varphi''(t) \, dt = \int_0^1 (1-t) \left( \sum_{i=1}^n \sum_{j=1}^n \frac{\partial^2 f}{\partial x_i \partial x_j}(\mathbf{x}_0 + t\mathbf{h}) \, h_i \, h_j \right) dt.
\]

\textbf{Εκτίμηση του υπολοίπου:}

Αφού η \(f\) είναι \(C^2\) στο \(D\), όλες οι δεύτερες μερικές παράγωγοι \(\frac{\partial^2 f}{\partial x_i \partial x_j}\) είναι συνεχείς στο \(D\). Επομένως, είναι φραγμένες στο συμπαγές σύνολο \(K\). Για κάθε \(i, j \in \{1, \dots, n\}\), ορίζουμε:
\[
M_{ij} = \max_{\mathbf{x} \in K} \left| \frac{\partial^2 f}{\partial x_i \partial x_j}(\mathbf{x}) \right| < \infty.
\]

Για κάθε \(t \in [0,1]\), το σημείο \(\mathbf{x}_0 + t\mathbf{h}\) ανήκει στο ευθύγραμμο τμήμα μεταξύ \(\mathbf{x}_0\) και \(\mathbf{x}_0 + \mathbf{h}\). Εφόσον το τμήμα αυτό περιέχεται εξ υποθέσεως στο \(K\), έχουμε \(\mathbf{x}_0 + t\mathbf{h} \in K\) για όλα τα \(t\).

Συνεπώς, για κάθε \(t \in [0,1]\):
\begin{align*}
|\varphi''(t)| &= \left| \sum_{i=1}^n \sum_{j=1}^n \frac{\partial^2 f}{\partial x_i \partial x_j}(\mathbf{x}_0 + t\mathbf{h}) \, h_i \, h_j \right| \\
&\le \sum_{i=1}^n \sum_{j=1}^n \left| \frac{\partial^2 f}{\partial x_i \partial x_j}(\mathbf{x}_0 + t\mathbf{h}) \right| \, |h_i| \, |h_j| \\
&\le \sum_{i=1}^n \sum_{j=1}^n M_{ij} \, |h_i| \, |h_j|.
\end{align*}

Θέτουμε \(M_{\max} = \max_{1 \le i,j \le n} M_{ij}\). Τότε:
\begin{align*}
|\varphi''(t)| &\le M_{\max} \sum_{i=1}^n \sum_{j=1}^n |h_i| \, |h_j| \\
&= M_{\max} \left( \sum_{i=1}^n |h_i| \right)^2.
\end{align*}

Χρησιμοποιούμε την ανισότητα Cauchy–Schwarz ή την ισοδυναμία νορμών στον \(\mathbb{R}^n\):
\[
\sum_{i=1}^n |h_i| \le \sqrt{n} \, \|\mathbf{h}\|.
\]
(Αυτό προκύπτει από το γεγονός ότι \(\sum |h_i| = \langle (|h_1|,\dots,|h_n|), (1,\dots,1) \rangle \le \|\mathbf{h}\| \cdot \|\mathbf{1}\| = \sqrt{n} \|\mathbf{h}\|\).)

Άρα:
\[
\left( \sum_{i=1}^n |h_i| \right)^2 \le n \|\mathbf{h}\|^2,
\]
και συνεπώς:
\[
|\varphi''(t)| \le M_{\max} \cdot n \cdot \|\mathbf{h}\|^2.
\]

Θέτουμε \(M_K' = n \cdot M_{\max}\). Τότε \(|\varphi''(t)| \le M_K' \|\mathbf{h}\|^2\) για όλα τα \(t \in [0,1]\).

Το υπόλοιπο φράσσεται τώρα ως εξής:
\begin{align*}
|R_1(\mathbf{h})| &= \left| \int_0^1 (1-t) \, \varphi''(t) \, dt \right| \\
&\le \int_0^1 (1-t) \, |\varphi''(t)| \, dt \\
&\le M_K' \|\mathbf{h}\|^2 \int_0^1 (1-t) \, dt \\
&= M_K' \|\mathbf{h}\|^2 \cdot \left[ t - \frac{t^2}{2} \right]_0^1 \\
&= M_K' \|\mathbf{h}\|^2 \cdot \frac{1}{2}.
\end{align*}

Ορίζοντας \(M_K = M_K' / 2 = \frac{n}{2} M_{\max}\), λαμβάνουμε το ζητούμενο φράγμα:
\[
|R_1(\mathbf{h})| \le M_K \cdot \|\mathbf{h}\|^2.
\]

Αυτό ολοκληρώνει την απόδειξη.
\end{proof}

\begin{remark}[Ύπαρξη του μεγίστου \(M_{ij}\)]
Η ύπαρξη των μεγίστων \(M_{ij}\) είναι εγγυημένη διότι κάθε \(\frac{\partial^2 f}{\partial x_i \partial x_j}\) είναι συνεχής συνάρτηση στο συμπαγές σύνολο \(K\). Από το θεώρημα ακροτάτων τιμών του Weierstrass, μια συνεχής συνάρτηση σε συμπαγές σύνολο επιτυγχάνει μέγιστη και ελάχιστη τιμή. Επομένως τα \(M_{ij}\) είναι πεπερασμένα και καλά ορισμένα.
\end{remark}

\begin{remark}[Εναλλακτική σταθερά μέσω νόρμας Εσσιανού]
Μια πιο κομψή διατύπωση χρησιμοποιεί τη φασματική νόρμα του Εσσιανού πίνακα. Αν ορίσουμε
\[
\|H_f(\mathbf{x})\|_2 = \sup_{\|\mathbf{v}\|=1} |\mathbf{v}^T H_f(\mathbf{x}) \mathbf{v}|,
\]
και θέσουμε \(M_K' = \max_{\mathbf{x} \in K} \|H_f(\mathbf{x})\|_2\), τότε:
\[
|\varphi''(t)| = |\mathbf{h}^T H_f(\mathbf{x}_0 + t\mathbf{h}) \mathbf{h}| \le M_K' \|\mathbf{h}\|^2,
\]
και η απόδειξη ολοκληρώνεται με \(M_K = M_K' / 2\). Αυτή η προσέγγιση αποφεύγει τον παράγοντα \(n\) και είναι βέλτιστη ως προς τη σταθερά.
\end{remark}

\begin{corollary}[Μετάφραση σε συμβολισμό \(O\)]
\label{cor:O_translation_Rn}
Στο πλαίσιο διαμέρισης όπου \(h = \max_i |h_i|\) (ή \(h = \operatorname{diam}(Q)\)), έχουμε \(\|\mathbf{h}\| \le \sqrt{n} \, h\). Το Λήμμα \ref{lem:C2_Taylor_bound_Rn} δίνει:
\[
|R_1(\mathbf{h})| \le M_K \|\mathbf{h}\|^2 \le M_K \cdot n \cdot h^2 = C_K \cdot h^2,
\]
όπου \(C_K = n M_K\). Συνεπώς:
\[
R_1(\mathbf{h}) = O(\|\mathbf{h}\|^2) \quad \text{και} \quad R_1(\mathbf{h}) = O(h^2).
\]
\end{corollary}

\begin{remark}[Γενίκευση σε \(C^k\) συναρτήσεις]
Το Λήμμα γενικεύεται άμεσα για \(C^{k+1}\) συναρτήσεις: το υπόλοιπο \(k\)-τάξης ικανοποιεί \(R_k(\mathbf{h}) = O(\|\mathbf{h}\|^{k+1})\). Η απόδειξη χρησιμοποιεί το ολοκληρωτικό υπόλοιπο \(k\)-τάξης:
\[
R_k(\mathbf{h}) = \int_0^1 \frac{(1-t)^k}{k!} \varphi^{(k+1)}(t) \, dt,
\]
και φράσσεται με τον ίδιο τρόπο από τα μέγιστα των παραγώγων τάξης \(k+1\) στο \(K\).
\end{remark}

\begin{remark}[Πρακτική εφαρμογή στα επιφανειακά ολοκληρώματα]
Στο Παράρτημα Α, η συνάρτηση \(f(u,v) = \|\mathbf{r}_u \times \mathbf{r}_v\|\) είναι \(C^2\) (εφόσον η \(\mathbf{r}\) είναι \(C^2\) και κανονική). Το Λήμμα \ref{lem:C2_Taylor_bound_Rn} με \(n=2\) εφαρμόζεται απευθείας για να δώσει \(R_1(\delta u, \delta v) = O(\delta u^2 + \delta v^2) = O(h^2)\), το οποίο είναι κρίσιμο για την τοπική εκτίμηση σφάλματος.
\end{remark}


\section{Θεώρημα μεταφοράς Reynolds}

TODO

\backmatter

\begin{thebibliography}{9}
\addcontentsline{toc}{chapter}{Bibliography}
% TODO: add references
\end{thebibliography}

\end{document}
